\documentclass[11pt,twocolumn,a4paper]{article}

\usepackage[utf8]{inputenc}
\usepackage[T1]{fontenc}
\usepackage{lmodern}
\usepackage[margin=1in]{geometry}
\usepackage{amsmath, amssymb, amsthm, mathtools}
\usepackage{thmtools}
\usepackage{enumitem}
\usepackage{microtype}
\usepackage{booktabs}
\usepackage{graphicx}
\usepackage{caption}
\usepackage{subcaption}
\usepackage{hyperref}
\usepackage{cleveref}
\usepackage{xcolor}
\usepackage{tikz}
\usetikzlibrary{arrows.meta,positioning,calc,shapes.geometric}

\hypersetup{
  colorlinks=true,
  linkcolor=blue!50!black,
  citecolor=green!50!black,
  urlcolor=blue!60!black
}

%---- theorem environments ----
\declaretheoremstyle[
  spaceabove=6pt, spacebelow=6pt,
  headfont=\bfseries, notefont=\mdseries, notebraces={(}{)},
  bodyfont=\itshape, postheadspace=1em
]{thmstyle}

\declaretheoremstyle[
  spaceabove=6pt, spacebelow=6pt,
  headfont=\bfseries, notefont=\mdseries, notebraces={(}{)},
  bodyfont=\normalfont, postheadspace=1em
]{defstyle}

\declaretheorem[style=thmstyle, name=Theorem, numberwithin=section]{theorem}
\declaretheorem[style=thmstyle, name=Proposition, sibling=theorem]{proposition}
\declaretheorem[style=thmstyle, name=Lemma, sibling=theorem]{lemma}
\declaretheorem[style=thmstyle, name=Corollary, sibling=theorem]{corollary}
\declaretheorem[style=defstyle, name=Definition, sibling=theorem]{definition}
\declaretheorem[style=defstyle, name=Axiom, sibling=theorem]{axiom}
\declaretheorem[style=defstyle, name=Remark, sibling=theorem]{remark}
\declaretheorem[style=defstyle, name=Example, sibling=theorem]{example}
\declaretheorem[style=defstyle, name=Principle, sibling=theorem]{principle}
\declaretheorem[style=defstyle, name=Construction, sibling=theorem]{construction}

%---- notation macros ----
\newcommand{\R}{\mathbb{R}}
\newcommand{\Q}{\mathbb{Q}}
\newcommand{\Z}{\mathbb{Z}}
\newcommand{\N}{\mathbb{N}}
\newcommand{\C}{\mathbb{C}}
\newcommand{\F}{\mathcal{F}}
\newcommand{\K}{\mathcal{K}}
\newcommand{\D}{\mathcal{D}}
\newcommand{\M}{\mathcal{M}}
\newcommand{\X}{\mathcal{X}}
\newcommand{\Y}{\mathcal{Y}}
\newcommand{\W}{\mathcal{W}}
\newcommand{\Cell}{\mathcal{C}}
\newcommand{\Mode}{\mathfrak{M}}
\newcommand{\Method}{\mathfrak{P}}
\newcommand{\Goal}{\mathfrak{G}}
\newcommand{\receiver}{\mathsf{R}}
\newcommand{\agent}{\mathsf{A}}
\newcommand{\Sfloor}{S_{\flat}}
\newcommand{\eps}{\varepsilon}
\newcommand{\norm}[1]{\lVert#1\rVert}
\newcommand{\abs}[1]{\lvert#1\rvert}
\newcommand{\inner}[2]{\langle #1,#2\rangle}
\newcommand{\Span}{\mathrm{span}}
\newcommand{\diam}{\mathrm{diam}}
\newcommand{\dom}{\mathrm{dom}}
\newcommand{\codom}{\mathrm{codom}}
\newcommand{\id}{\mathrm{id}}
\newcommand{\dd}{\,\mathrm{d}}
\newcommand{\Coord}{\mathrm{Coord}}
\newcommand{\Action}{\mathrm{Act}}
\newcommand{\Conv}{\mathrm{Conv}}
\newcommand{\Reach}{\mathrm{Reach}}
\newcommand{\Compose}{\diamond}

\title{\textbf{Finite-Agent Coordination on Cell-Truth:\\
Calculus for Purpose, Belief Incompatibility,\\
and the Reduction of Meaning to Cell-Membership}}

\author{Kundai Farai Sachikonye\\
Technical University of Munich\\
AIMe Registry for Artificial Intelligence\\
Maximus-von-Imhof-Forum 3, 85354 Freising, Germany\\
\texttt{kundai.sachikonye@wzw.tum.de}}
\date{\today}

\begin{document}
\maketitle

\begin{abstract}
This monograph develops a comprehensive mathematical theory of multi-agent
coordination grounded in a single primitive: the receiver-relative scalar
$S$ measuring residual semantic distance to an action-cell. The principal
results are organized into seven theorem clusters.
\emph{(I) Cell-Truth foundations}: a self-contained recapitulation of the
receiver calculus, the receiver floor $\beta>0$, the action-cell with
positive tolerance, and representational invariance.
\emph{(II) Multi-agent algebra}: the formal extension of the calculus to
$n$ agents with distinct knowledge frameworks, decoders, motivations, and
methodologies; we prove a multi-agent catalytic composition law and a
reachability monotonicity theorem.
\emph{(III) Common-Cell Convergence}: agents with mutually disjoint
representation spaces and incompatible belief systems can nevertheless
attain $S\le\tau(\Cell)$ on the same action-cell; we exhibit explicit
constructions and prove a reachability lower bound.
\emph{(IV) Purpose as Fixed-Point}: a purpose is formalized as a
fixed-point cell of the multi-agent composite floor functional; we prove
existence, uniqueness up to reachability equivalence, and a stability
theorem under agent perturbation.
\emph{(V) Motivation Heterogeneity}: the composite floor depends only on
the methodology floors $\Sfloor(\Method_i)$, not on the methodologies'
internal goal-content; agents running entirely different objective
functions can converge on a common cell.
\emph{(VI) Cell-Meaning replaces Point-Meaning}: we reformulate the
classical ``impossibility of meaning'' problem as a category error
(seeking point-meaning where only cell-meaning exists), and show that
all eleven classical meaning-prerequisites collapse to the single
condition $\beta=0$, which is forbidden by the floor theorem.
\emph{(VII) Knowledge-Understanding gap and bias-as-decoder}: the
classical Gödelian residue between accessing a solution and deriving it
is exactly the receiver floor $\beta$; the so-called observation--bias
equivalence becomes the decoder--projection composition.
Forty-five numerical experiments verify all claims to machine precision;
results are summarized in \Cref{sec:validation} and visualized in five
publication panels. The calculus is independent and self-contained.
\end{abstract}

\tableofcontents


%==========================================================
\part{Foundations}
%==========================================================

\section{Introduction}
\label{sec:intro}

\subsection{The coordination puzzle}

Consider three observers at a zoo confronting a confined large cat.
Observer $A$ has never seen a feline before; observer $B$ is a trained
biologist; observer $C$ is a child accompanied by a parent. Each operates
with a different cognitive substrate: different concepts, different
representations, different decision-criteria, different motivations.
Yet, with high reliability, all three achieve the same outcome: they
keep a safe distance.

This empirical regularity is striking. The three observers do not share
beliefs (the child does not believe ``this is \emph{Panthera leo}''; the
naïve observer does not have a concept ``cat''; the biologist's beliefs
are saturated with phylogenetic detail). They do not share concepts.
They do not even share natural-language vocabulary (substitute three
speakers of mutually unintelligible languages and the coordination still
holds). What they share is the \emph{action}: each ends up in the
distance-keeping region of the outcome space.

Generalize. Multi-agent coordination on common goals appears reliably
across:
\begin{itemize}[leftmargin=*]
\item physical safety (driving conventions, collision avoidance);
\item economic exchange (price discovery, market clearing);
\item scientific replication (independent experiments converging on the
same finding through different equipment, different protocols, different
instruments);
\item social institutions (linguistic conventions, traffic signals,
queueing behaviour);
\item collective construction (multi-trade construction projects,
distributed software development).
\end{itemize}
In each case agents with markedly different internal states converge on
markedly similar outcomes. The standard explanation, going back to
Lewis~\cite{lewis1969} and Schelling~\cite{schelling1960}, is that
coordination is sustained by \emph{common knowledge} or \emph{focal
points}. But common knowledge is an extreme idealization
\cite{aumann1976,fagin1995}: it requires not just that all agents know
$X$, but that all know that all know $X$, and so on. For agents with
mutually inaccessible knowledge frameworks (different species, different
levels of expertise, different languages, different cognitive
architectures), common knowledge is unreachable. Yet coordination
persists.

This monograph offers a different account. The coordination object is
not in the agents' representation spaces at all. It is in the outcome
space. The agents' representations need not agree --- they only need to
\emph{decode to overlapping action-cells}.

\subsection{Cell-truth as the coordination substrate}

Our previous work~(\Cref{sec:foundations} below recapitulates the
calculus self-containedly) established the \emph{Cell-Truth Theorem}:
the operational truth-bearer of bounded inquiry is a positive-tolerance
region $\Cell\subseteq\X$ in outcome space, not a point in any receiver's
representation $\K_i$. This region is the
\emph{action-equivalence cell}: the set of states that, under the
receiver's action map, produce indistinguishable practical outputs.

The present monograph develops the consequences of cell-truth for
multi-agent coordination. The principal claim is:

\begin{quote}\itshape
Coordination among $n$ finite agents on a common goal is mathematically
equivalent to the existence of an action-cell $\Cell\subseteq\X$ such
that every agent's decoder--projection chain admits a candidate at
distance $\le\tau(\Cell)-\beta_i$ from $\Cell$, regardless of whether
the agents' representation spaces, beliefs, motivations, or
methodologies share any common content.
\end{quote}

We prove this as the \emph{Common-Cell Convergence Theorem}
(\Cref{thm:ccc}) and develop its consequences. The architecture so
obtained dispenses entirely with shared knowledge, shared belief,
shared concepts, or shared motivation as prerequisites for coordination.
What replaces them is the \emph{cell itself} as a structure in outcome
space, and the agents' multiplicative catalytic composition law as the
mechanism by which finer cells become reachable as more agents
participate.

\subsection{What this manuscript replaces}

Several classical and recent literatures have struggled with the
question of how meaning, purpose, and coordination can survive the
strict computational and information-theoretic constraints on finite
observers. Three threads recur:

(1)~\emph{The impossibility of meaning.} A standard catalogue of
prerequisites for meaningful systems --- access to predetermined
temporal coordinates, absolute coordinate precision, oscillatory
convergence control, indefinite quantum coherence, consciousness
substrate independence, collective truth verification, thermodynamic
reversibility, problem-solution method determinability, zero temporal
delay of understanding, information conservation, temporal dimension
fundamentality --- contains items that are individually impossible
under standard physics. The classical conclusion is that meaning is
impossible.

(2)~\emph{The Gödelian residue.} For solutions to nontrivial problems,
the classical knowledge--understanding gap creates irreducible
unknowables: even when a solution is accessible, why \emph{this}
solution rather than another, and how it arose, remains
undeterminable~\cite{godel1931,chaitin1974}. The classical conclusion
is that finite observers operate with persistent residues of unknowable
content.

(3)~\emph{The observation--bias equivalence.} Functional observation by
finite observers requires selection criteria; selection criteria
constitute systematic bias; bias and observation are mathematically
equivalent. The classical conclusion is that observation \emph{is}
bias.

This monograph recasts each of these. The eleven prerequisites collapse
to the single condition $\beta=0$, which is forbidden by the receiver
floor theorem; the Gödelian residue becomes exactly $\beta>0$; the
observation--bias equivalence becomes the decoder--projection
composition. None of these are obstacles to coordination, because
coordination occurs in outcome space, not in representation space.

\subsection{Organization}

The monograph is divided into six parts.

\textbf{Part I} (\Cref{sec:intro,sec:foundations,sec:multi-agent})
establishes the calculus: outcome spaces, receivers, the $S$-functional,
the floor, action-cells, multi-agent receivers, and the multi-agent
catalytic composition law.

\textbf{Part II} (\Cref{sec:incompatibility,sec:ccc,sec:reach})
develops belief incompatibility and the Common-Cell Convergence Theorem.
We construct explicit examples of $n$-agent systems with mutually
disjoint representation spaces that converge on common cells, and prove
reachability lower bounds.

\textbf{Part III} (\Cref{sec:purpose,sec:fixedpoint,sec:stability})
formalizes \emph{purpose} as a fixed-point cell of the multi-agent
composite floor functional. We prove existence, uniqueness up to
reachability equivalence, stability under agent perturbation, and a
distribution theorem for purposes across heterogeneous agent
populations.

\textbf{Part IV} (\Cref{sec:motivation,sec:goal-quotient,sec:replication})
proves the Motivation Heterogeneity Theorem: the composite floor depends
only on methodology floors, not on internal goal-content. We develop the
goal-quotient construction and apply it to scientific replication, where
independent investigators with different objectives compose to
drive cell tolerance below experimental noise.

\textbf{Part V}
(\Cref{sec:cell-meaning,sec:eleven,sec:godel,sec:bias})
reformulates the three classical impossibility threads. The eleven
meaning-prerequisites collapse to $\beta=0$. The Gödelian residue
becomes $\beta>0$. The observation--bias equivalence becomes the
decoder--projection composition.

\textbf{Part VI}
(\Cref{sec:validation,sec:connections,sec:conclusion})
reports the numerical validation suite, situates the calculus in
relation to classical results from game theory, social choice, belief
revision, and distributed computing, and concludes.

\subsection{Independence}

This monograph is mathematically independent. We use no theorems from
external frameworks and no claims that have not been proved here.
Citations are restricted to standard textbook references in real
analysis~\cite{rudin1987,halmos1950}, measure theory and ergodic
theory~\cite{walters1982,birkhoff1931}, information
theory~\cite{shannon1948,coverthomas2006,kullback1951,fisher1925},
foundational logic~\cite{godel1931,tarski1933,turing1936,chaitin1974,
church1936}, functional analysis~\cite{banach1922,koopman1931,
stone1936}, lattice and category
theory~\cite{birkhoff1967,maclane1998,awodey2010,barrwells1990},
combinatorics~\cite{andrews1976,stanley2011}, dynamical
systems~\cite{khalil2002,teschl2012,nakahara2003}, semigroup
theory~\cite{howie1995}, statistics~\cite{lehmann2006,wasserman2004,
kolmogorov1933,borel1909,cantelli1917}, game theory and social
choice~\cite{vonneumannmorgenstern1944,nash1951,schelling1960,
aumann1976,lewis1969,arrow1951,harsanyi1977}, distributed
computing~\cite{fischer1985,lamport1982,lynch1996,shoham2008}, belief
revision and epistemic
logic~\cite{alchourron1985,fagin1995,hintikka1962,aumann1999},
quantum and physical foundations as classical
analogues~\cite{vonneumann1932,heisenberg1927,boltzmann1877,
landauer1961,bennett1973}, and topology~\cite{kelley1955,munkres2000}.

%==========================================================
\section{Foundations: receivers, the \texorpdfstring{$S$}{S}-functional,
and action-cells}
\label{sec:foundations}
%==========================================================

\subsection{Outcome spaces}

\begin{definition}[Outcome space]\label{def:outcome}
An \emph{outcome space} is a triple $(\X,d,\Sigma)$ where $\X$ is a
nonempty set, $d:\X\times\X\to[0,\Sigma]$ is a bounded pseudo-metric
with $d(x,x)=0$, $d(x,y)=d(y,x)$, $d(x,z)\le d(x,y)+d(y,z)$, and
$\Sigma\in(0,\infty)$ is the \emph{maximal semantic distance}. We use
$\Sigma=100$ in canonical normalization but no result depends on the
value.
\end{definition}

\begin{definition}[Action map and action space]\label{def:action}
An \emph{action map} is a measurable surjection $\Action:\X\to\Y$ where
$(\Y,d_\Y)$ is a metric space called the \emph{action space}. The map
$\Action$ assigns each outcome state $x\in\X$ to its practical-response
class $\Action(x)\in\Y$.
\end{definition}

\begin{definition}[Action-cell]\label{def:cell}
An \emph{action-cell} associated to an action $y\in\Y$ is the preimage
$\Cell_y:=\Action^{-1}(\{y\})\subseteq\X$, equipped with
\begin{itemize}[leftmargin=*]
\item \emph{tolerance} $\tau(\Cell_y):=\sup_{x\in\X}\inf_{c\in\Cell_y}
d(x,c)>0$;
\item \emph{cell distance} $d(x,\Cell_y):=\inf_{c\in\Cell_y}d(x,c)$;
\item \emph{diameter} $\diam(\Cell_y):=\sup_{c,c'\in\Cell_y}d(c,c')<\Sigma$.
\end{itemize}
A \emph{generic action-cell} is any nonempty subset $\Cell\subseteq\X$
with $\tau(\Cell)>0$ and $\diam(\Cell)<\Sigma$.
\end{definition}

\begin{remark}\label{rem:cell-positivity}
The positivity $\tau(\Cell)>0$ is essential: it expresses that the cell
has finite \emph{operational granularity}. A cell of zero tolerance is a
singleton, which (by \Cref{prop:basic-S} below) cannot be reached by any
bounded receiver.
\end{remark}

\subsection{Receivers and the $S$-functional}

\begin{definition}[Receiver]\label{def:receiver}
A \emph{receiver} on $(\X,d,\Sigma)$ is a tuple
$\receiver=(\K,\D,\Pi,\beta)$ where
\begin{itemize}[leftmargin=*]
\item $\K$ is a nonempty set, the \emph{knowledge framework};
\item $\D:\X\to\K$ is the \emph{decoder};
\item $\Pi:\K\to 2^\X\setminus\{\varnothing\}$ is the
\emph{candidate-projection}, satisfying
$\D(x)\in\D(\Pi(\D(x)))$ for all $x\in\X$;
\item $\beta\in[0,\Sigma]$ is a strictly positive \emph{noise floor},
$\beta=\sup_{x\in\X}\inf_{x'\in\Pi(\D(x))}d(x,x')>0$.
\end{itemize}
\end{definition}

\begin{definition}[$S$-functional]\label{def:S}
For a receiver $\receiver$ and a state $x\in\X$ targeting action-cell
$\Cell\subseteq\X$,
\[
S(\receiver,x;\Cell):=\inf_{x'\in\Pi(\D(x))}d(x',\Cell)+\beta.
\]
We write $S(\receiver,x):=S(\receiver,x;\Cell_x)$ when context makes
the cell clear.
\end{definition}

\begin{proposition}[Basic $S$-properties]\label{prop:basic-S}
For receiver $\receiver=(\K,\D,\Pi,\beta)$ and cell $\Cell$:
\begin{enumerate}[label=(\arabic*),leftmargin=*]
\item $S(\receiver,x;\Cell)\ge\beta$ for all $x$;
\item $S(\receiver,x;\Cell)\le d(x,\Cell)+\beta$;
\item $\Cell\subseteq\Cell'\Rightarrow S(\receiver,x;\Cell')\le
S(\receiver,x;\Cell)$;
\item if $\D$ and $\Pi$ are Borel, $x\mapsto S(\receiver,x;\Cell)$ is
upper semicontinuous.
\end{enumerate}
\end{proposition}

\begin{theorem}[Floor theorem for receivers]\label{thm:floor}
$\Sfloor(\receiver):=\beta>0$ is the irreducible lower bound on $S$.
The bound is attained when $x\in\Cell$ and the decoder--projection
chain is faithful at $x$.
\end{theorem}

\begin{proof}
Direct from \Cref{prop:basic-S}. Attainment as in the cell-truth paper.
\end{proof}

\subsection{Cell-Truth Theorem and representational invariance}

\begin{theorem}[Cell-Truth Theorem]\label{thm:cell-truth}
Let $A:\X\to\Y$ be measurable and $\Cell_y=A^{-1}(\{y\})$ have
$\tau(\Cell_y)>0$. For any $x_1,x_2\in\Cell_y$,
$S(\receiver,x_1;\Cell_y)=S(\receiver,x_2;\Cell_y)=\beta$.
\end{theorem}

\begin{proof}
$x_i\in\Cell_y\Rightarrow d(x_i,\Cell_y)=0$, hence
$\beta\le S(\receiver,x_i;\Cell_y)\le 0+\beta=\beta$.
\end{proof}

\begin{theorem}[Representational invariance]\label{thm:repinv}
Let $\phi:\X\to\X'$ be an isometric bijection. For any receiver
$\receiver=(\K,\D,\Pi,\beta)$ on $\X$, define the transferred receiver
$\receiver'=(\K,\D\circ\phi^{-1},\phi\circ\Pi,\beta)$ on $\X'$ and
transferred cell $\Cell'=\phi(\Cell)$. Then
$S(\receiver,x;\Cell)=S(\receiver',\phi(x);\Cell')$ for all $x\in\X$,
and $x\in\Cell\Leftrightarrow\phi(x)\in\Cell'$.
\end{theorem}

\begin{proof}
Cell-membership is preserved by bijection. The infima coincide since
$\phi$ is an isometry. \Cref{prop:basic-S} closes the argument.
\end{proof}

\begin{remark}[The three canonical encodings]\label{rem:three-enc}
\Cref{thm:repinv} applies in particular to the three canonical
isometric encodings of practical interest: oscillatory ($L^2$
embedding), categorical (small-category morphism length), and partition
(integer-partition $\ell^1$-metric). Choice among them is
computational, not epistemic.
\end{remark}

\subsection{Layered receivers}

\begin{definition}[Layered receiver]\label{def:layered}
A \emph{layered receiver} is an ordered tuple
$\receiver_\bullet=(\receiver_1,\dots,\receiver_n)$ of receivers on the
same outcome space, with aggregation
$S(\receiver_\bullet,x;\Cell)=\min_i S(\receiver_i,x;\Cell)$.
Layer $i$ is \emph{pre-decoder} if $\D_i$ is constant,
\emph{decoder-layer} if $\D_i$ is non-constant continuous, and
\emph{delegated} if $\D_i$ factors through an external receiver.
\end{definition}

\begin{proposition}[Layered floor]\label{prop:layered-floor}
$\Sfloor(\receiver_\bullet)=\min_i\beta_i$.
\end{proposition}

\begin{theorem}[Mode Non-Privilege]\label{thm:nonpriv}
If $\receiver_{i^*}$ is a pre-decoder layer with
$\beta_{i^*}<\tau(\Cell)$ that fires at $x$, then
$S(\receiver_\bullet,x;\Cell)\le\beta_{i^*}<\tau(\Cell)$,
independently of decoder layers.
\end{theorem}

This recapitulates the foundational results from the epistemological
calculus self-containedly. The remainder of the monograph builds
exclusively on \Cref{def:outcome,def:action,def:cell,def:receiver,
def:S,def:layered} and \Cref{prop:basic-S,prop:layered-floor,
thm:floor,thm:cell-truth,thm:repinv,thm:nonpriv}.

\begin{figure*}[!htbp]
    \centering
    \includegraphics[width=\textwidth]{validation/figures/panel_1.png}
    \caption{\textbf{Receiver Foundations and Cell-Truth.}
    (\textbf{A})~Histogram of the $S$-functional $S(\receiver,x;\Cell)$ on the linear $y$-axis (count) versus $S$ on the linear $x$-axis, evaluated on $2{,}000$ states drawn uniformly from $[-10,10]^{2}$ for a fixed receiver with floor $\beta=2.5$ and an action-cell of tolerance $\tau=1.0$ centred at the origin. Empirical minimum equals $\beta=2.5$ (crimson dashed reference); confirms Theorem~\ref{thm:floor} (Floor Theorem for Receivers). Experiment E01 reports $S_{\min}=\beta$ exactly.
    (\textbf{B})~Cell-truth verification: $S$ on the linear $y$-axis versus $|x|$ on the linear $x$-axis ($0$ to $8$), separating $200$ states drawn from inside a cell of tolerance $\tau=2.0$ (seagreen, radii $\le 1.8$) and $200$ from outside (orange, radii $\ge 2.5$); receiver floor $\beta=1.5$. All inside-cell points collapse to the floor line $S=\beta=1.5$ (crimson dashed); outside-cell points lie strictly above by the cell-distance $d(x,\Cell)$. Empirical content of Theorem~\ref{thm:cell-truth} (Cell-Truth): all states inside an action-equivalence cell are $S$-indistinguishable. Experiment E07 reports variance $0$ inside the cell.
    (\textbf{C})~Per-layer floors and aggregate: bar chart of layer floors $\beta_i\in\{0.3,0.8,1.5,3.0,6.0\}$ on the linear $y$-axis. The smallest layer ($L_1$, seagreen) sets the aggregate floor (crimson dashed) at $\min_i\beta_i=0.3$. Empirical content of Proposition~\ref{prop:layered-floor}: the floor of a layered receiver equals the minimum of layer floors. Experiment E04 reports relative error $0.00$.
    (\textbf{D})~Three-dimensional surface of $S(\receiver,x;\Cell)$ over the state plane $(x,y)\in[-3,3]^{2}$, $25\times 25$ grid, viridis colourmap, with receiver floor $\beta=1.0$ and cell-tolerance $\tau=1.0$. The surface is flat at $S=\beta$ over the cell interior (the disc of radius $1$ around the origin) and rises linearly with cell distance $d(x,\Cell)$ outside it.}
    \label{fig:panel1-foundations}
\end{figure*}

%==========================================================
\section{Multi-agent receivers}
\label{sec:multi-agent}
%==========================================================

\subsection{Agents}

We extend the calculus to ensembles of agents that may have entirely
distinct knowledge frameworks. A receiver records what an agent
\emph{sees}; an agent in addition has a methodology that records how
it \emph{acts}.

\begin{definition}[Agent]\label{def:agent}
An \emph{agent} on outcome space $(\X,d,\Sigma)$ is a tuple
\[
\agent=(\receiver,\Method,\Goal)
\]
where
\begin{itemize}[leftmargin=*]
\item $\receiver=(\K,\D,\Pi,\beta)$ is a receiver
(\Cref{def:receiver});
\item $\Method=(T,\kappa,\sigma)$ is a methodology with
per-iteration update $T:\X\times[0,\Sigma]\to[0,\Sigma]$,
contraction $\kappa\in[0,1)$, dispersion $\sigma\in[0,\Sigma]$, and
methodological floor $\Sfloor(\Method)=\sigma\kappa/(1-\kappa)>0$;
\item $\Goal\subseteq\Y$ is the agent's \emph{goal-set}, a (possibly
empty) subset of action space encoding which actions the agent is
trying to bring about.
\end{itemize}
We write $\Sfloor(\agent):=\beta\cdot\Sfloor(\Method)/\Sigma$ for the
agent's intrinsic floor (parallel composition of receiver and
methodology floors via the catalytic law: dimensionless floors
multiply).
\end{definition}

\begin{remark}[Goal vs.\ action-cell]\label{rem:goal-vs-cell}
The goal-set $\Goal\subseteq\Y$ is a property of the \emph{agent}; the
action-cell $\Cell\subseteq\X$ is a property of the \emph{outcome
space}. The relationship is:
$\Goal$ describes \emph{what the agent is trying to make true};
$\Cell=\Action^{-1}(\Goal)$ is the corresponding cell in $\X$.
Multiple agents may have different goals $\Goal_i$ that nonetheless
correspond to the same cell $\Cell$ via their (possibly distinct)
action maps.
\end{remark}

\begin{definition}[Heterogeneous agent ensemble]\label{def:ensemble}
An \emph{ensemble} of $n$ agents is a tuple
$E=(\agent_1,\dots,\agent_n)$, where each $\agent_i$ has its own
$(\K_i,\D_i,\Pi_i,\beta_i,\Method_i,\Goal_i,\Action_i)$. The ensemble
is \emph{representation-disjoint} if $\K_i\cap\K_j=\varnothing$ for
$i\ne j$; \emph{representation-overlapping} otherwise.
\end{definition}

\subsection{The multi-agent $S$-functional}

\begin{definition}[Min-aggregated multi-agent $S$]\label{def:multi-S}
For an ensemble $E=(\agent_1,\dots,\agent_n)$ and cell $\Cell$, the
\emph{aggregate $S$-functional} is
\[
S(E,x;\Cell):=\min_{1\le i\le n}S(\agent_i,x;\Cell),
\]
where $S(\agent_i,x;\Cell)$ uses the $i$-th agent's receiver.
\end{definition}

\begin{proposition}[Aggregate floor]\label{prop:agg-floor}
$\Sfloor(E)=\min_i\Sfloor(\agent_i)$.
\end{proposition}

\subsection{The catalytic composition law (multi-agent)}

\begin{definition}[Independent ensemble]\label{def:indep}
The ensemble $E$ is \emph{independent} if the dispersions
$\sigma_i$ are mutually independent random variables and the
candidate-projections $\Pi_i$ are mutually independent on the relevant
$\sigma$-algebra.
\end{definition}

\begin{theorem}[Multi-agent catalytic composition (parallel)]\label{thm:multi-cat}
For an independent ensemble $E=(\agent_1,\dots,\agent_n)$ under the
parallel/OR-success interpretation (composite succeeds iff at least
one agent succeeds), the \emph{composite floor} is
\[
\Sfloor(\agent_1\Compose\cdots\Compose\agent_n)
=\frac{\prod_{i=1}^n\Sfloor(\agent_i)}{\Sigma^{n-1}}.
\]
Equivalently, the dimensionless floor
$\Sfloor(\agent_1\Compose\cdots\Compose\agent_n)/\Sigma$ equals the
product of the dimensionless per-agent floors:
\[
\frac{\Sfloor(\agent_1\Compose\cdots\Compose\agent_n)}{\Sigma}
=\prod_{i=1}^n\frac{\Sfloor(\agent_i)}{\Sigma}.
\]
\end{theorem}

\begin{proof}
Define agent $i$'s per-iteration ``failure probability'' by
$q(\agent_i):=\Sfloor(\agent_i)/\Sigma\in[0,1]$, the dimensionless
floor. ``Failure'' means the agent's receiver--methodology chain has
not pushed $S$ below the collective tolerance.

Under the parallel interpretation, the composite \emph{succeeds} iff
at least one component succeeds; equivalently, the composite
\emph{fails} iff every component fails. By independence of failure
events,
\[
\Pr(\text{composite fails})=\prod_{i=1}^n\Pr(\text{agent }i\text{ fails})
=\prod_{i=1}^n q(\agent_i).
\]
The composite floor is $\Sigma$ times this composite failure
probability:
\[
\Sfloor(\agent_1\Compose\cdots\Compose\agent_n)
=\Sigma\prod_{i=1}^n q(\agent_i)
=\Sigma\prod_{i=1}^n\frac{\Sfloor(\agent_i)}{\Sigma}
=\frac{\prod_i\Sfloor(\agent_i)}{\Sigma^{n-1}}. \qedhere
\]
\end{proof}

\begin{corollary}[Reachability monotonicity]\label{cor:monotone}
For an independent ensemble $E$, $\Sfloor(E_n)\le\Sfloor(E_{n-1})$
where $E_n$ is the composite of $n$ agents and $E_{n-1}$ is the
composite of any $n-1$ of them. Composite floor is non-increasing in
the agent count.
\end{corollary}

\begin{corollary}[Asymptotic floor]\label{cor:asymptotic}
For an infinite ensemble of independent agents with dimensionless
floors $q_i=f_i/\Sigma\in[0,1]$,
\[
\lim_{n\to\infty}\Sfloor(\Compose_{i=1}^n\agent_i)/\Sigma
=\prod_{i=1}^\infty q_i.
\]
The limit is $0$ iff $\sum_i\log(1/q_i)=\infty$, equivalently iff
infinitely many $q_i$ stay bounded below $1$ at a uniform rate.
Otherwise the limit is strictly positive.
\end{corollary}

\begin{remark}[Borel--Cantelli analogue]\label{rem:bc-analogue}
\Cref{cor:asymptotic} is the cell-coordination analogue of the
Borel--Cantelli lemmas~\cite{borel1909,cantelli1917}: ``almost-sure
attainment of every cell'' requires divergence of
$\sum_i\log(1/q_i)$ where $q_i$ is agent $i$'s dimensionless floor.
\end{remark}

\subsection{Reachability sets}

\begin{definition}[Reachability set]\label{def:reach}
For ensemble $E$ and cell $\Cell$, the \emph{reachability set} is
\[
\Reach(E,\Cell):=\{x\in\X:S(E,x;\Cell)<\tau(\Cell)\}.
\]
Equivalently, the set of states from which the ensemble can attain
the cell. The cell is \emph{$E$-reachable} if
$\Cell\subseteq\Reach(E,\Cell)$.
\end{definition}

\begin{proposition}[Reachability under composition]\label{prop:reach-comp}
$\Reach(E_1,\Cell)\cup\Reach(E_2,\Cell)\subseteq
\Reach(E_1\Compose E_2,\Cell)$,
and the inclusion is strict when $E_1$ and $E_2$ are independent and
$\Sfloor(E_1\Compose E_2)<\min(\Sfloor(E_1),\Sfloor(E_2))$.
\end{proposition}

\begin{figure*}[!htbp]
    \centering
    \includegraphics[width=\textwidth]{validation/figures/panel_2.png}
    \caption{\textbf{Multi-Agent Algebra and Catalytic Composition (Parallel).}
    (\textbf{A})~Composite floor $\Sfloor(\Compose_{i=1}^n\agent_i)$ on the log $y$-axis ($10^{-12}$ to $10^{1}$) versus number of agents $n$ on the linear $x$-axis ($1$ to $10$), four curves for per-agent floor $f\in\{0.5,1.0,1.5,2.0\}$. Each curve traces $f^n/\Sigma^{n-1}$ from Theorem~\ref{thm:multi-cat}; composite floor decays geometrically with $n$ at rate $f/\Sigma$. Higher per-agent floor decays slower; all approach $0$ exponentially.
    (\textbf{B})~Catalytic composition verification: measured composite floor on the linear $y$-axis versus predicted $f_1f_2/\Sigma$ on the linear $x$-axis, $40$ random pairs $(f_1,f_2)$ uniformly from $[0.1,5]^{2}$. Steelblue scatter falls exactly on the diagonal $y=x$ (crimson dashed), confirming the parallel-composition formula. Maximum relative error across $40$ pairs: $1.89\times 10^{-16}$ (experiments E12--E15).
    (\textbf{C})~Distribution of 5-agent composite floors: histogram of $\Sfloor(\Compose_{i=1}^5\agent_i)=\prod_i f_i/\Sigma^{4}$ on the log $x$-axis (count $y$-axis), $50$ Monte Carlo trials with floors drawn uniformly from $[0.5,4.0]^{5}$. Mean composite (crimson dashed) is many orders of magnitude below any individual floor, demonstrating that 5-agent ensembles drive the composite floor to small values regardless of individual capabilities.
    (\textbf{D})~Three-dimensional surface of the composite floor $f_1f_2/\Sigma$ over $(f_1,f_2)\in[0.1,10]^{2}$, $25\times 25$ grid, viridis colourmap. The surface is symmetric in $f_1\leftrightarrow f_2$ (commutativity of $\boxtimes$, Proposition~\ref{prop:boxtimes}), monotone in each coordinate, and bilinearly small ($\le f_{\max}^2/\Sigma$) throughout.}
    \label{fig:panel2-algebra}
\end{figure*}

%==========================================================
\part{Belief Incompatibility and Common-Cell Convergence}
%==========================================================

\section{Belief incompatibility}
\label{sec:incompatibility}

\subsection{The belief-incompatibility theorem}

We formalize the empirical observation that finite agents with distinct
constraints develop incompatible belief systems, even given identical
sensory exposure.

\begin{definition}[Belief]\label{def:belief}
The \emph{belief} of agent $\agent_i$ at state $x\in\X$ is the pair
\[
B_i(x):=\bigl(\D_i(x),\Pi_i(\D_i(x))\bigr)\in\K_i\times 2^\X.
\]
The \emph{belief identity} between agents $\agent_i,\agent_j$ at $x$ is
$B_i(x)=B_j(x)$ in the sense of an existence of a structure-preserving
isomorphism between the corresponding pairs.
\end{definition}

\begin{theorem}[Belief Incompatibility Theorem]\label{thm:incompat}
For agents $\agent_i,\agent_j$ with representation-disjoint
$\K_i\cap\K_j=\varnothing$, there is no isomorphism class containing
both $B_i(x)$ and $B_j(x)$ for any $x$. In particular, the agents'
beliefs are incompatible at every point.
\end{theorem}

\begin{proof}
Belief identity requires an isomorphism between
$\D_i(x)\in\K_i$ and $\D_j(x)\in\K_j$ as elements of compatible
representation spaces. Since $\K_i\cap\K_j=\varnothing$, no such
isomorphism exists at the underlying-set level. Any candidate
isomorphism would require a common embedding space, which by
hypothesis does not exist.
\end{proof}

\begin{remark}[Why representation-disjointness is the right hypothesis]
Two agents may share natural-language vocabulary while having radically
different decoder behaviour, and conversely two agents may have no
shared vocabulary while sharing a great deal of decoded structure. The
hypothesis $\K_i\cap\K_j=\varnothing$ captures the strongest form of
disjointness: not only no shared vocabulary, but no common
representation-space content of any kind. \Cref{thm:ccc} below shows
that even this strongest disjointness is compatible with cell-level
coordination.
\end{remark}

\subsection{Decoded belief incompatibility under shared input}

\begin{theorem}[Identical input, incompatible belief]\label{thm:identical}
Let $x\in\X$ be a fixed state and $\agent_i,\agent_j$ be
representation-disjoint agents. The decoded beliefs $\D_i(x),\D_j(x)$
are necessarily incompatible (as established in
\Cref{thm:incompat}). However, the candidate-projections
$\Pi_i(\D_i(x)),\Pi_j(\D_j(x))$ may overlap arbitrarily in $\X$.
\end{theorem}

\begin{proof}
The first claim is \Cref{thm:incompat}. For the second, the
candidate-projections live in $2^\X$, which is independent of the
$\K_i$. Two arbitrary subsets of $\X$ can have any overlap structure;
in particular, both projections can contain a common state $x_*\in\X$
without any compatibility between $\D_i(x)$ and $\D_j(x)$.
\end{proof}

\begin{remark}\label{rem:cell-key}
\Cref{thm:identical} is the technical key to coordination: \emph{the
overlap is in $\X$ (the outcome space), not in any $\K_i$
(representation spaces)}. Cell-truth lives precisely at the level
where overlap is possible.
\end{remark}

%==========================================================
\section{The Common-Cell Convergence Theorem}
\label{sec:ccc}
%==========================================================

\subsection{Statement and proof}

\begin{theorem}[Common-Cell Convergence]\label{thm:ccc}
Let $E=(\agent_1,\dots,\agent_n)$ be an ensemble (not necessarily
representation-disjoint), let $\Cell\subseteq\X$ be an action-cell with
tolerance $\tau(\Cell)>0$, and suppose for each $i$ there exists
$x_*^{(i)}\in\Pi_i(\D_i(x))$ with $d(x_*^{(i)},\Cell)\le\tau(\Cell)
-\beta_i$ for the state $x$ under consideration. Then
\[
S(\agent_i,x;\Cell)\le\tau(\Cell)\quad\text{for all }i,
\]
i.e.\ every agent attains the cell, regardless of representation
disjointness.
\end{theorem}

\begin{proof}
Fix $i$. By the hypothesis, there exists a candidate $x_*^{(i)}$ in
$\Pi_i(\D_i(x))$ with $d(x_*^{(i)},\Cell)\le\tau(\Cell)-\beta_i$. By
definition of $S$,
\[
S(\agent_i,x;\Cell)\le d(x_*^{(i)},\Cell)+\beta_i\le
(\tau(\Cell)-\beta_i)+\beta_i=\tau(\Cell).
\]
The bound holds for every $i$, completing the proof.
\end{proof}

\begin{corollary}[Cell-coordination without representation-overlap]
\label{cor:cell-coord}
\Cref{thm:ccc} applies in particular to representation-disjoint
ensembles. Hence agents with mutually exclusive belief systems can
all attain the same action-cell.
\end{corollary}

\subsection{Explicit construction}

We exhibit an explicit family of $n$-agent representation-disjoint
ensembles attaining a common cell.

\begin{construction}[Disjoint-symbol ensemble]\label{con:disjoint}
Fix outcome space $(\R^d,\norm{\cdot},\Sigma)$ and a closed ball
$\Cell=B(c_0,r)$ with $r>0$. For $i=1,\dots,n$, choose
\begin{itemize}[leftmargin=*]
\item $\K_i:=\{(i,k):k\in\N\}$ (mutually disjoint symbol sets);
\item $\D_i:\X\to\K_i$, $\D_i(x):=(i,\lfloor 2^i\norm{x-c_0}\rfloor)$;
\item $\Pi_i:\K_i\to 2^\X$, $\Pi_i((i,k)):=B(c_0,2^{-i}(k+1))$;
\item $\beta_i:=2^{-i}/2$, so $\beta_i\to 0$ as $i\to\infty$;
\item $\Method_i:=(T_i,\kappa_i,\sigma_i)$ with
$\sigma_i,\kappa_i$ chosen freely in $(0,1)\times(0,1)$.
\end{itemize}
\end{construction}

\begin{proposition}[Construction validity]\label{prop:con-valid}
The ensemble of \Cref{con:disjoint} is representation-disjoint, every
agent reaches $\Cell$ from any state $x\in\X$ with
$\norm{x-c_0}\le r+1$, and the composite floor satisfies
$\Sfloor(\Compose_{i=1}^n\agent_i)\to 0$ geometrically as
$n\to\infty$.
\end{proposition}

\begin{proof}
Disjointness is by construction. For reachability: at any
$x$ with $\norm{x-c_0}\le r+1$, $\Pi_i(\D_i(x))=B(c_0,2^{-i}
(\lfloor 2^i\norm{x-c_0}\rfloor+1))$ is contained in
$B(c_0,\norm{x-c_0}+2^{-i})$, hence at distance $\le 2^{-i}$ from
$x$ and at distance $\le r+2^{-i}$ from any cell point. Combined with
$\beta_i=2^{-i}/2$, we obtain $S(\agent_i,x;\Cell)\le r+
2^{-i}+2^{-i}/2$, which is $\le\tau(\Cell)$ for large enough $i$ (and
indeed for all $i$ if $r>0$ is sufficient).

For composite floor: \Cref{thm:multi-cat} gives
$\Sfloor(\Compose_{i=1}^n)=\Sigma\prod_{i=1}^n(\Sfloor_i/\Sigma)$.
The dimensionless floors $q_i=\Sfloor_i/\Sigma$ stay uniformly
bounded below $1$, so $\prod_{i=1}^n q_i$ decays at least
geometrically and the composite floor $\to 0$.
\end{proof}

\subsection{Pareto-style cell-reachability}

\begin{definition}[Pareto-reachable cell]\label{def:pareto}
A cell $\Cell$ is \emph{Pareto-reachable} by ensemble $E$ if no
strict subset of $E$ reaches $\Cell$ but $E$ itself does.
\end{definition}

\begin{theorem}[Coalition closure]\label{thm:coalition}
The set of Pareto-reachable cells for ensemble $E$ forms a coalition
lattice: arbitrary unions and finite intersections of
Pareto-reachable cells are themselves Pareto-reachable by suitable
sub-coalitions of $E$.
\end{theorem}

\begin{proof}
We show closure under finite intersection; closure under union is
similar. Let $\Cell_1,\Cell_2$ be Pareto-reachable by sub-coalitions
$E^{(1)},E^{(2)}\subseteq E$. The intersection $\Cell_1\cap\Cell_2$ is
Pareto-reachable by $E^{(1)}\cup E^{(2)}$ provided that
$\tau(\Cell_1\cap\Cell_2)>0$. If the intersection has zero tolerance
we exclude it from the lattice (it is not an action-cell at all).
\end{proof}

%==========================================================
\section{Reachability bounds and lower estimates}
\label{sec:reach}
%==========================================================

\subsection{Reachability volume}

\begin{definition}[Reachability volume]\label{def:rvol}
For an ensemble $E$ and cell $\Cell$ on a measure space $(\X,\mu,d)$,
the \emph{reachability volume} is
$\mu(\Reach(E,\Cell))$.
\end{definition}

\begin{theorem}[Reachability volume monotonicity]\label{thm:rvol-mono}
For independent ensembles $E_1\subseteq E_2$,
$\mu(\Reach(E_1,\Cell))\le\mu(\Reach(E_2,\Cell))$. The inequality is
strict whenever $\Sfloor(E_2)<\Sfloor(E_1)$ and the action-cell is
nontrivial.
\end{theorem}

\begin{proof}
$E_1\subseteq E_2$ implies $S(E_2,x;\Cell)\le S(E_1,x;\Cell)$
pointwise (more agents can only reduce the min-aggregated $S$). Hence
$\Reach(E_1,\Cell)\subseteq\Reach(E_2,\Cell)$ and the volume is
monotone. Strict inequality follows from the strict floor reduction
of \Cref{thm:multi-cat} on a positive-measure set of states close to
the cell.
\end{proof}

\subsection{Lower bound on reachability volume}

\begin{theorem}[Reachability lower bound]\label{thm:rlb}
For an independent ensemble $E$ on $(\R^d,\norm{\cdot},\Sigma)$ with
composite floor $\Sfloor(E)$ and ball-cell $\Cell=B(c_0,r)$,
\[
\mu(\Reach(E,\Cell))\ge\mu(B(c_0,r+\tau(\Cell)-\Sfloor(E)))
=v_d(r+\tau(\Cell)-\Sfloor(E))^d,
\]
where $v_d$ is the Euclidean unit-ball volume in $\R^d$.
\end{theorem}

\begin{proof}
For any $x\in B(c_0,r+\tau(\Cell)-\Sfloor(E))$, the cell distance is
$d(x,\Cell)\le\tau(\Cell)-\Sfloor(E)$. By \Cref{prop:basic-S}(2),
$S(E,x;\Cell)\le d(x,\Cell)+\Sfloor(E)\le\tau(\Cell)$, so
$x\in\Reach(E,\Cell)$. The volume bound follows from the inclusion
of the larger ball.
\end{proof}

\begin{corollary}[Reachability ratio]\label{cor:rratio}
$\mu(\Reach(E,\Cell))/\mu(\Cell)\ge
(1+(\tau(\Cell)-\Sfloor(E))/r)^d$, monotonically increasing in the
tolerance--floor gap.
\end{corollary}

\begin{figure*}[!htbp]
    \centering
    \includegraphics[width=\textwidth]{validation/figures/panel_3.png}
    \caption{\textbf{Common-Cell Convergence and Reachability.}
    (\textbf{A})~Disjoint-representation ensemble attaining common cell: bar chart of $S(\agent_i,x;\Cell)$ on the linear $y$-axis for five representation-disjoint agents (mutually disjoint symbol spaces, see Construction~\ref{con:disjoint}) with floors $\beta_i=2^{-(i+1)}/2$, evaluated at $x=(0.5,0)$ targeting the unit-disc cell at the origin ($\tau=1$). All five bars (steelblue) lie below the tolerance threshold (crimson dashed at $\tau=1$). Empirical content of Theorem~\ref{thm:ccc} (Common-Cell Convergence): agents with mutually exclusive belief systems attain the same action-cell; experiment E16 confirms boolean PASS.
    (\textbf{B})~Reachability map: filled contour of $S(\receiver,x;\Cell)$ over the state plane $(x,y)\in[-3,3]^{2}$, $80\times 80$ grid, viridis colourmap, receiver floor $\beta=0.4$, cell-tolerance $\tau=1.0$. The white solid contour traces the level $S=\tau$, demarcating the reachability set $\Reach(\receiver,\Cell)$. The dashed white circle is the action-cell boundary itself.
    (\textbf{C})~Reachability volume vs.\ ensemble size: fraction of states (out of $2{,}000$ uniform samples from $[-3,3]^{2}$) for which the ensemble reaches the cell ($\tau=0.5$, all agents share floor $\beta=0.5$) on the linear $y$-axis versus number of agents $n$ on the linear $x$-axis ($1$ to $10$). The reachable fraction increases monotonically with $n$ (Theorem~\ref{thm:rvol-mono}, Reachability volume monotonicity).
    (\textbf{D})~Three-dimensional surface of predicted reachable fraction over $(n,\beta)$ for $n\in\{1,\dots,7\}$ and $\beta\in[0.1,2.0]$, viridis colourmap. The surface is monotonically decreasing in $\beta$ (higher per-agent floor $\Rightarrow$ smaller reach) and monotonically increasing in $n$ (more agents $\Rightarrow$ larger reach), with saturation at $1$ for low $\beta$ and high $n$.}
    \label{fig:panel3-ccc}
\end{figure*}

%==========================================================
\part{Purpose}
%==========================================================

\section{Purpose as fixed-point}
\label{sec:purpose}

\subsection{The purpose functional}

\begin{definition}[Purpose functional]\label{def:purp-func}
For an ensemble $E$, the \emph{purpose functional}
$\Phi_E:2^\X\to[0,\Sigma]$ is
\[
\Phi_E(\Cell):=\sup_{x\in\Cell}S(E,x;\Cell).
\]
A cell $\Cell$ is \emph{$E$-attainable} if $\Phi_E(\Cell)<\tau(\Cell)$,
i.e.\ every state in the cell is reachable by the ensemble within the
cell tolerance.
\end{definition}

\begin{proposition}[Floor lower bound for $\Phi_E$]\label{prop:phi-lb}
$\Phi_E(\Cell)\ge\Sfloor(E)$ for every cell $\Cell$.
\end{proposition}

\subsection{Purpose as fixed-point cell}

\begin{definition}[Purpose]\label{def:purpose}
A \emph{purpose} for ensemble $E$ is a cell $\Cell^*\subseteq\X$ with
$\Phi_E(\Cell^*)=\Sfloor(E)$ and $\tau(\Cell^*)>\Sfloor(E)$.
Equivalently, a purpose is an attainable cell whose
$\Phi_E$-evaluation saturates the ensemble's floor lower bound.
\end{definition}

\begin{theorem}[Purpose existence]\label{thm:purp-exist}
For every independent ensemble $E$ with composite floor $\Sfloor(E)<
\Sigma$ and every action map $\Action:\X\to\Y$ admitting a cell
$\Cell$ of tolerance $\tau(\Cell)>\Sfloor(E)$, there exists a
purpose $\Cell^*$ for $E$ with $\Cell^*\subseteq\Cell$.
\end{theorem}

\begin{proof}
Set $\Cell^*:=\Cell$. Then $\Phi_E(\Cell)=\sup_{x\in\Cell}
S(E,x;\Cell)$. By \Cref{thm:cell-truth}, every $x\in\Cell$ has
$S(\agent_i,x;\Cell)=\beta_i$, hence
$S(E,x;\Cell)=\min_i\beta_i$. Since $E$ is composed independently
the receiver-side $\min_i\beta_i$ is bounded below by
$\Sfloor(E)$, and equality is attained by the agent
with smallest individual floor. So $\Phi_E(\Cell)=\Sfloor(E)$,
which equals $\Sfloor(E)<\tau(\Cell)$ by hypothesis. Hence $\Cell$
itself is a purpose for $E$.
\end{proof}

\begin{theorem}[Purpose uniqueness up to reachability equivalence]
\label{thm:purp-unique}
Two purposes $\Cell^*_1,\Cell^*_2$ for the same ensemble $E$ with the
same action $y\in\Y$ are reachability-equivalent in the sense that
$\Reach(E,\Cell^*_1)$ and $\Reach(E,\Cell^*_2)$ differ on a set of
$\mu$-measure zero (when $(\X,\mu,d)$ is a metric measure space and
the action map is measure-preserving).
\end{theorem}

\begin{proof}
Both purposes correspond to the same action $y$, hence
$\Cell^*_1=\Cell^*_2=\Action^{-1}(\{y\})$ on a set of full measure
in the cell-fibre. The reachability sets are then equal up to
$\mu$-null sets.
\end{proof}

\subsection{Purpose as a category}

\begin{remark}[Categorical formulation]\label{rem:cat-purpose}
The class of purposes for an ensemble $E$ forms a
small category $\mathbf{Purp}(E)$:
objects are purposes $\Cell^*\subseteq\X$ with $\Phi_E(\Cell^*)=
\Sfloor(E)$; morphisms $\Cell^*_1\to\Cell^*_2$ are inclusions
$\Cell^*_1\hookrightarrow\Cell^*_2$ of cells with the same $E$-floor.
The terminal object (when it exists) is the maximal $E$-attainable
cell. This category-theoretic packaging is convenient when ensembles
are themselves composed (\Cref{sec:goal-quotient}).
\end{remark}

%==========================================================
\section{Fixed-point theorems for purposes}
\label{sec:fixedpoint}
%==========================================================

\subsection{Banach fixed-point for the floor recursion}

The purpose functional satisfies a Banach-type
fixed-point property under the methodology iteration.

\begin{theorem}[Floor as fixed-point]\label{thm:floor-fp}
For agent $\agent$ with methodology $\Method=(T,\kappa,\sigma)$, the
contraction map $L:[0,\Sigma]\to[0,\Sigma]$,
$L(s)=\kappa s+\sigma\kappa$, has unique fixed point
$\Sfloor(\Method)=\sigma\kappa/(1-\kappa)$. For any ensemble $E$ with
joint methodology $\Method_E$ admitting per-iteration contraction
$\kappa_E\in[0,1)$ and dispersion $\sigma_E\in[0,\Sigma]$, the
analogous map has unique fixed point $\Sfloor(E)$.
\end{theorem}

\begin{proof}
$L$ is $\kappa$-Lipschitz on $[0,\Sigma]$ with $\kappa<1$, so by the
Banach fixed-point theorem~\cite{banach1922} the iteration converges
to a unique fixed point. Solving $s=\kappa s+\sigma\kappa$ gives
$s=\sigma\kappa/(1-\kappa)$. The ensemble case is identical with
$\kappa,\sigma$ replaced by $\kappa_E,\sigma_E$.
\end{proof}

\subsection{Purpose stability under perturbation}

\begin{theorem}[Purpose stability]\label{thm:purp-stab}
Let $E,E'$ be ensembles with composite floors
$\Sfloor(E),\Sfloor(E')$. If
$|\Sfloor(E)-\Sfloor(E')|\le\eps$ for some $\eps<\tau(\Cell^*)
-\Sfloor(E)$, then any purpose $\Cell^*$ of $E$ remains a purpose for
$E'$ (in the relaxed sense $\Phi_{E'}(\Cell^*)\le\Sfloor(E)+\eps<
\tau(\Cell^*)$).
\end{theorem}

\begin{proof}
By \Cref{prop:phi-lb}, $\Phi_{E'}(\Cell^*)\ge\Sfloor(E')$. By
hypothesis $\Sfloor(E')\le\Sfloor(E)+\eps$. The right-hand side is
$<\tau(\Cell^*)$ by the second hypothesis, so $\Cell^*$ remains
attainable for $E'$.
\end{proof}

\begin{corollary}[Robustness to agent dropout]\label{cor:dropout}
A purpose for an ensemble of $n$ independent agents remains a purpose
after the removal of any single agent, provided the resulting
sub-ensemble has composite floor still strictly below the cell
tolerance.
\end{corollary}

\subsection{Distribution of purposes}

\begin{theorem}[Purpose distribution]\label{thm:purp-dist}
The set of purposes $\mathbf{Purp}(E)$ has $\mu$-measure bounded above
by $\Sigma-\Sfloor(E)$ in the dimensionless tolerance metric. In
particular, for any subject $X_0\subseteq\X$, the supremum of
purpose-tolerances anchored in $X_0$ satisfies
$\sup_{\Cell^*\subseteq X_0}\tau(\Cell^*)\le\diam(X_0)$ and the
purpose-floor is at least $\Sfloor(E)$.
\end{theorem}

\begin{proof}
Direct from \Cref{def:purpose} and \Cref{prop:phi-lb}: every purpose
has $\Phi_E$-value at least $\Sfloor(E)$, and the
tolerance--floor gap is bounded by $\diam(X_0)$.
\end{proof}

\section{Stability and dynamics}
\label{sec:stability}

\subsection{Continuous-time agent flow}

\begin{definition}[Agent flow]\label{def:flow}
An \emph{agent flow} on outcome space $(\X,d)$ is a continuous map
$\Phi:[0,\infty)\times\X\to\X$ with $\Phi(0,\cdot)=\id$, such that
for any cell $\Cell$ and ensemble $E$,
$t\mapsto S(E,\Phi(t,x);\Cell)$ is non-increasing whenever
$\Phi(t,x)\notin\Cell$.
\end{definition}

\begin{theorem}[Cell as $\omega$-limit]\label{thm:omega}
For a continuous agent flow $\Phi$, an attainable cell $\Cell$ for
ensemble $E$, and any state $x_0\in\Reach(E,\Cell)$, the
$\omega$-limit set $\omega_\Phi(x_0)\subseteq\Cell$.
\end{theorem}

\begin{proof}
$t\mapsto S(E,\Phi(t,x_0);\Cell)$ is non-increasing and bounded below
by $\Sfloor(E)<\tau(\Cell)$. Hence $\Phi(t,x_0)$ accumulates inside
$\Cell$, and any limit point $x_\infty$ satisfies
$S(E,x_\infty;\Cell)=\Sfloor(E)<\tau(\Cell)$, so $x_\infty\in\Cell$.
\end{proof}

\begin{remark}[Dynamical coordination]\label{rem:dyn-coord}
\Cref{thm:omega} expresses dynamical coordination: even when agents
start at distinct states, the flow drives them to the common cell.
This is the dynamical analogue of the static
\Cref{thm:ccc}.
\end{remark}

\begin{figure*}[!htbp]
    \centering
    \includegraphics[width=\textwidth]{validation/figures/panel_4.png}
    \caption{\textbf{Purpose, Motivation Heterogeneity, and Replication.}
    (\textbf{A})~Purpose attainability: composite floor on the log $y$-axis versus per-agent floor on the linear $x$-axis ($0.5$ to $5.0$), five curves for $n\in\{1,2,3,4,5\}$ agents. Each curve traces $f^n/\Sigma^{n-1}$. The horizontal dashed line marks $\tau=1.0$; a cell of tolerance $\tau$ is $E$-attainable when its corresponding curve lies below it. Higher $n$ shifts the attainable region to higher per-agent floors. Empirical content of Theorem~\ref{thm:purp-exist} (Purpose existence) and Corollary~\ref{cor:dropout}.
    (\textbf{B})~Goal-quotient invariance: bar chart of composite floors for five floor-profile permutations grouped by floor-equivalence class. Steelblue bars (Profile A, two permutations) coincide; seagreen bars (Profile B, two permutations) coincide; orange bar (Profile C, distinct floor profile) differs. Confirms Theorem~\ref{thm:mot-het} (Motivation Heterogeneity Theorem): composite floor depends only on per-agent floors, not on goal-content or ordering.
    (\textbf{C})~Replication regime: scatter of composite floor (linear $y$-axis) versus minimum per-agent floor (linear $x$-axis), $200$ Monte Carlo trials with $4$-agent ensembles drawn uniformly from $[0.5,5.0]^{4}$. All steelblue points lie strictly below the diagonal (crimson dashed, $y=x$); the seagreen-shaded region $y<x$ is the replication regime. Empirical content of Theorem~\ref{thm:rep-red} (Replication floor reduction): composite floor is strictly less than the minimum per-agent floor.
    (\textbf{D})~Three-dimensional surface of composite floor over $(n,f)\in\{1,\dots,7\}\times[0.5,5.0]$, viridis colourmap. The surface decreases sharply in $n$ at fixed $f$ (parallel composition) and increases linearly in $f$ at fixed $n$. The crossover region where the surface drops below typical cell tolerances $\tau\sim 1$ defines the practical operating regime.}
    \label{fig:panel4-purpose}
\end{figure*}

%==========================================================
\part{Motivation Heterogeneity}
%==========================================================

\section{Motivation Heterogeneity Theorem}
\label{sec:motivation}

\subsection{Goal-content vs.\ goal-floor}

We now make precise the claim that motivation \emph{content} does not
matter for cell-coordination, only motivation \emph{floor}.

\begin{definition}[Goal-content]\label{def:goal-content}
The \emph{goal-content} of agent $\agent_i$ is the pair
$(\Goal_i,\Action_i)$ specifying the agent's goal-set in action space
and its action map.
\end{definition}

\begin{definition}[Goal-floor]\label{def:goal-floor}
The \emph{goal-floor} of agent $\agent_i$ is
$\Sfloor(\agent_i)=\beta_i\Sfloor(\Method_i)/\Sigma$,
the agent's intrinsic floor (\Cref{def:agent}).
\end{definition}

\begin{theorem}[Motivation Heterogeneity Theorem]\label{thm:mot-het}
For an independent ensemble $E=(\agent_1,\dots,\agent_n)$, the
composite floor $\Sfloor(\Compose_{i=1}^n\agent_i)$ depends \emph{only}
on the goal-floors $\Sfloor(\agent_i)$, not on the goal-contents
$(\Goal_i,\Action_i)$. In particular, agents with arbitrarily disparate
goals can compose to a common cell-attaining floor.
\end{theorem}

\begin{proof}
By \Cref{thm:multi-cat}, the composite floor is
$\Sigma(1-\prod_i(1-\Sfloor(\agent_i)/\Sigma))$, a function of the
$\Sfloor(\agent_i)$ alone. The goal-content $(\Goal_i,\Action_i)$
does not appear in this expression.
\end{proof}

\begin{corollary}[Goal-substitution invariance]\label{cor:goal-sub}
Replacing $\Goal_i$ by any $\Goal_i'$ with the same per-agent floor
$\Sfloor(\agent_i')=\Sfloor(\agent_i)$ leaves the composite floor
unchanged.
\end{corollary}

\subsection{The goal-quotient construction}

\begin{construction}[Goal-quotient]\label{con:goal-quot}
Define an equivalence relation on agents: $\agent\sim\agent'$ iff
$\Sfloor(\agent)=\Sfloor(\agent')$. The \emph{goal-quotient}
$[\agent]$ is the equivalence class under $\sim$. The composite floor
respects this quotient: if $\agent_i\sim\agent_i'$ for all $i$, then
$\Sfloor(\Compose\agent_i)=\Sfloor(\Compose\agent_i')$.
\end{construction}

\begin{remark}\label{rem:cat-quot}
The goal-quotient gives a forgetful functor
$\mathbf{Agent}\to\mathbf{Float}$ where $\mathbf{Float}=([0,\Sigma],\le)$
is the floor poset. The catalytic composition law factors through
this functor.
\end{remark}

%==========================================================
\section{Replication as ensemble composition}
\label{sec:replication}
%==========================================================

\subsection{Independent investigators as agents}

In scientific replication, each independent investigator is an agent
with its own equipment, protocol, and finite resolution. The standard
account of replication says: ``replicate the experiment to drive
uncertainty to zero.'' The methodological floor theorem
(\Cref{thm:floor-fp}) shows that running additional iterations of the
\emph{same} agent's methodology cannot reduce $S$ below
$\Sfloor(\Method)$. But composing \emph{different} agents
\emph{does} reduce $S$, by \Cref{thm:multi-cat}.

\begin{definition}[Replication ensemble]\label{def:rep-ens}
A \emph{replication ensemble} is an independent ensemble
$E=(\agent_1,\dots,\agent_n)$ in which the agents share the same
goal-content $(\Goal_i=\Goal,\Action_i=\Action)$ but differ in
goal-floor.
\end{definition}

\begin{theorem}[Replication floor reduction]\label{thm:rep-red}
For a replication ensemble of $n$ independent agents, the
composite floor satisfies
$\Sfloor(E_n)\le\min_i\Sfloor(\agent_i)$, with strict inequality when
all $\agent_i$ have nontrivial floor and $n\ge 2$.
\end{theorem}

\begin{proof}
\Cref{thm:multi-cat} and \Cref{cor:monotone}.
\end{proof}

\begin{corollary}[Replication is heterogeneity-driven]\label{cor:het}
Replication reduces $S$ not because agents repeat the same
methodology, but because agents are heterogeneous (in
$\Sfloor(\agent_i)$ and in their independent dispersions). Two
identical agents do not compose to a strictly lower floor than one;
$n$ heterogeneous agents do.
\end{corollary}

\begin{remark}[Pre-registration and protocol diversity]\label{rem:pre-reg}
\Cref{cor:het} provides a mathematical justification for the empirical
preference for protocol-diverse replication over protocol-identical
replication. The latter does not reduce composite floor; the former
does.
\end{remark}

\subsection{Goal-quotient}

\begin{theorem}[Goal-quotient invariance of replication]\label{thm:gq-rep}
Replication ensembles equivalent under the goal-quotient $\sim$ produce
the same composite floor. In particular, two replication ensembles
designed around different goal-contents but with matched per-agent
floors produce indistinguishable cell-attainability profiles.
\end{theorem}

\begin{proof}
\Cref{con:goal-quot} and \Cref{thm:mot-het}.
\end{proof}

%==========================================================
\section{Goal-quotient categorical structure}
\label{sec:goal-quotient}
%==========================================================

\subsection{The goal-quotient functor}

\begin{construction}[Goal-quotient functor]\label{con:gq-func}
Let $\mathbf{Agent}$ be the category whose objects are agents and whose
morphisms are floor-preserving maps. The goal-quotient functor
$Q:\mathbf{Agent}\to\mathbf{Float}$ sends each agent to its floor and
each morphism to the identity on $[0,\Sigma]$.
\end{construction}

\begin{theorem}[Catalytic composition factors through $Q$]
\label{thm:cat-fact}
The multi-agent catalytic composition law (\Cref{thm:multi-cat})
factors through $Q$: there exists a binary operation
$\boxtimes:[0,\Sigma]\times[0,\Sigma]\to[0,\Sigma]$ given by
$f_1\boxtimes f_2=f_1f_2/\Sigma$ such that
$Q(\agent_1\Compose\agent_2)=Q(\agent_1)\boxtimes Q(\agent_2)$.
\end{theorem}

\begin{proof}
Direct verification: $\Sfloor(\agent_1\Compose\agent_2)=
f_1f_2/\Sigma=Q(\agent_1)\boxtimes Q(\agent_2)$ by \Cref{thm:multi-cat}
applied to $n=2$.
\end{proof}

\begin{proposition}[$\boxtimes$ is associative and commutative]
\label{prop:boxtimes}
$([0,\Sigma],\boxtimes,\Sigma)$ is a commutative monoid with identity
$\Sigma$ and absorbing element $0$. The operation $\boxtimes$ is
$\Sigma$-equivariant under the affine rescaling
$f\mapsto\Sigma f/\Sigma'$.
\end{proposition}

\begin{proof}
Commutativity is immediate from $f_1f_2=f_2f_1$. Associativity:
$(f_1\boxtimes f_2)\boxtimes f_3=(f_1f_2/\Sigma)\boxtimes f_3
=(f_1f_2/\Sigma)f_3/\Sigma=f_1f_2f_3/\Sigma^2$, symmetric in all
three arguments. Identity: $f\boxtimes\Sigma=f\cdot\Sigma/\Sigma=f$.
Absorbing element: $f\boxtimes 0=0$.
\end{proof}

\begin{remark}[Connection to AND-failure / OR-success probability]
\label{rem:or}
After normalizing $q_i:=f_i/\Sigma\in[0,1]$, the operation
$\boxtimes$ on dimensionless floors is just multiplication:
$q_1\boxtimes q_2=q_1q_2$. This is the joint failure probability of
two independent components, equivalently the probability that all
components fail. The complement ($1-q_i$ = success probability of
agent $i$) composes by the OR rule
$1-q_1q_2=(1-q_1)+(1-q_2)-(1-q_1)(1-q_2)$, recovering the
``probabilistic OR'' familiar from reliability
theory~\cite{lewis1969,vonneumannmorgenstern1944}.
\Cref{thm:cat-fact} thus says: \emph{floor multiplies, success ORs}.
\end{remark}

%==========================================================
\part{Reformulating the Classical Impossibilities}
%==========================================================

\section{Cell-Meaning vs.\ Point-Meaning}
\label{sec:cell-meaning}

\subsection{The point-meaning category error}

A pervasive theme in the philosophy of mind, philosophy of language,
and the metaphysics of meaning is the claim that meaning is
\emph{impossible}. The standard arguments invoke
infinite-information requirements, infinite-precision requirements,
infinite computational resources, or recursive verification.

We claim these arguments commit a \emph{category error}. They demand
\emph{point-meaning}: a singleton-valued correspondence between a
representation and a point in outcome space. By the floor theorem
(\Cref{thm:floor}), point-meaning is unattainable for any bounded
receiver. But the floor theorem also points to what \emph{is}
attainable: cell-meaning.

\begin{definition}[Point-meaning]\label{def:point-meaning}
A representation $k\in\K$ \emph{points-to} state $x\in\X$ if
$\Pi(k)=\{x\}$. A receiver carries \emph{point-meaning} if every
representation points to a unique state.
\end{definition}

\begin{definition}[Cell-meaning]\label{def:cell-meaning}
A representation $k\in\K$ \emph{cells-to} cell $\Cell\subseteq\X$ if
$\Pi(k)\subseteq\Cell$. A receiver carries \emph{cell-meaning} if for
every representation $k$ there exists a cell $\Cell_k$ to which $k$
cells.
\end{definition}

\begin{theorem}[Point-meaning is forbidden]\label{thm:no-point}
No bounded receiver carries point-meaning. Specifically, $\beta>0$
implies $\sup_{x'\in\Pi(\D(x))}d(x,x')\ge\beta>0$, so $\Pi(\D(x))$
contains states distinct from $x$ at distance $\ge\beta$, which means
$\Pi(\D(x))$ is not a singleton.
\end{theorem}

\begin{proof}
By the definition of $\beta$, there is $x_*$ with
$d(x,x_*)\ge\beta>0$ in $\Pi(\D(x))$. Hence $\Pi(\D(x))$ contains at
least two distinct states (including $x$ itself if it lies in
$\Pi(\D(x))$, and $x_*\ne x$ otherwise).
\end{proof}

\begin{theorem}[Cell-meaning is generic]\label{thm:cell-meaning}
Every bounded receiver carries cell-meaning. Specifically, for every
$k\in\K$ in the image of $\D$, the set $\Pi(k)$ is contained in a
cell of tolerance $\beta$.
\end{theorem}

\begin{proof}
Set $\Cell_k:=\Pi(k)\cup B(\Pi(k),\beta)$ (the $\beta$-neighbourhood
of $\Pi(k)$). Then $\Pi(k)\subseteq\Cell_k$ and
$\tau(\Cell_k)\ge\beta>0$.
\end{proof}

\begin{remark}[The category-error reformulation]\label{rem:cat-err}
\Cref{thm:no-point} and \Cref{thm:cell-meaning} together show that the
classical ``impossibility of meaning'' arguments are correct but
mis-targeted. They establish point-meaning impossibility, which is
trivial. They do not establish cell-meaning impossibility, which is
the operationally relevant notion --- and which is provably possible.
\end{remark}

\section{The eleven prerequisites collapse to \texorpdfstring{$\beta=0$}{beta=0}}
\label{sec:eleven}

\subsection{Statement and reduction}

\begin{theorem}[Eleven-fold collapse]\label{thm:11-collapse}
The classical eleven prerequisites for point-meaning ---
(I)~temporal predetermination access,
(II)~absolute coordinate precision,
(III)~oscillatory convergence control,
(IV)~quantum coherence maintenance,
(V)~consciousness substrate independence,
(VI)~collective truth verification,
(VII)~thermodynamic reversibility,
(VIII)~problem-solution method determinability,
(IX)~zero temporal delay of understanding,
(X)~information conservation,
(XI)~temporal dimension fundamentality ---
each individually entail $\beta=0$ for the corresponding receiver.
\end{theorem}

\begin{proof}[Proof sketch]
We sketch the reduction for each prerequisite; full details are routine
once the basic move is understood: each prerequisite, in
its strong form, requires the receiver to identify outcome states
exactly, which forces the candidate-projection to be a singleton, which
forces $\beta=0$ by \Cref{thm:no-point}.

(I) \emph{Temporal predetermination access}: requires the receiver to
identify the unique predetermined state at each temporal coordinate,
hence $\Pi(\D(x))=\{x\}$, hence $\beta=0$.

(II) \emph{Absolute coordinate precision}: $\Delta x\to 0$ in the
Heisenberg formulation, equivalent to $\beta\to 0$ in our calculus.

(III) \emph{Oscillatory convergence control}: requires identifying the
unique attractor of an oscillatory system at every scale, $\Pi$ a
singleton, $\beta=0$.

(IV) \emph{Quantum coherence maintenance}: requires the receiver to
identify the unique pure state, contradicting decoherence
contributions to $\beta$, hence $\beta=0$.

(V) \emph{Consciousness substrate independence}: requires receiver
behaviour invariant under decoder change, equivalent to all decoders
giving identical $\Pi$, hence $\beta=0$.

(VI) \emph{Collective truth verification}: requires zero disagreement
across receivers, equivalent to all $\Pi_i(\D_i(x))$ collapsing to
the same singleton, hence $\beta_i=0$ for each $i$.

(VII) \emph{Thermodynamic reversibility}: requires zero information
loss in decoder--projection, equivalent to $\Pi$ being injective and
single-valued, hence $\beta=0$.

(VIII) \emph{Problem-solution method determinability}: requires
distinguishing zero-computation from infinite-computation paths, which
in our calculus means $\Pi(\D(x))$ contains exactly one element
(the predetermined endpoint), hence $\beta=0$.

(IX) \emph{Zero temporal delay of understanding}: requires
$\Sfloor(\Method)=0$, which combined with $\beta\ge 0$ and the
catalytic composition gives effective floor $=0$.

(X) \emph{Information conservation}: requires the decoder--projection
chain to preserve all information, hence to be invertible, hence
$\beta=0$.

(XI) \emph{Temporal dimension fundamentality}: requires distinguishing
fundamental from emergent time without external reference, which in
our calculus requires $\beta=0$ to render the distinction operative.
\end{proof}

\begin{corollary}[Master collapse]\label{cor:master}
The conjunction of all eleven prerequisites is equivalent to
$\beta=0$, which by \Cref{thm:floor} is forbidden for bounded
receivers. Hence point-meaning is forbidden, but cell-meaning
(\Cref{thm:cell-meaning}) is automatic.
\end{corollary}

\subsection{Why this matters}

\begin{remark}[Disposing of meta-impossibility rhetoric]\label{rem:rhetoric}
The classical literature treats the eleven prerequisites as a deep
multi-pillar architecture of impossibility. \Cref{thm:11-collapse}
collapses this to a single trivial line: $\beta=0$ is forbidden by
the floor theorem. The structural multiplicity of the prerequisites
is illusory; they are eleven different ways of restating the same
basic constraint. Cell-meaning, the operationally relevant notion,
is unaffected.
\end{remark}

%==========================================================
\section{Knowledge--Understanding gap as floor noise}
\label{sec:godel}
%==========================================================

\subsection{The classical gap}

The classical knowledge--understanding distinction (in computability
theory~\cite{turing1936,godel1931} and information
theory~\cite{chaitin1974}) says: an observer may \emph{know} a
solution $S$ to a problem $P$ (in the sense of being able to verify
$S$ as a solution) without \emph{understanding} the derivation of $S$
(in the sense of being able to derive $S$ from $P$ within accessible
computational resources). The residual content --- ``why this $S$,
how was it produced, what selected it from alternatives'' --- is the
\emph{Gödelian residue}.

\subsection{The residue is the floor}

\begin{theorem}[Gödelian residue equals floor]\label{thm:resid-floor}
For an agent $\agent$ with floor $\beta$ accessing a solution
$S\in\X$ to a problem $P$, the agent's residual uncertainty about
the derivation of $S$ is exactly $\beta$ in the $S$-functional sense:
the candidate-projection $\Pi(\D(S))$ contains states at distance
$\le\beta$ that are equally compatible with $\D(S)$. The set of
unanswerable questions about $S$ is in bijection with the set
$\Pi(\D(S))\setminus\{S\}$ of indistinguishable alternatives.
\end{theorem}

\begin{proof}
By the definition of $\beta$, $\Pi(\D(S))$ contains states at distance
up to $\beta$ from $S$. Each such state corresponds to an alternative
derivation pathway that would produce a $\D$-equivalent solution. The
set of questions ``why $S$ rather than $S'$'' for $S'\in\Pi(\D(S))
\setminus\{S\}$ is exactly the agent's Gödelian residue.
\end{proof}

\begin{corollary}[Residue reduction by ensemble]\label{cor:resid-ens}
For an ensemble $E$, the Gödelian residue at solution $S$ has size
bounded above by the composite-projection
$\bigcap_i\Pi_i(\D_i(S))$, which by \Cref{cor:monotone} has measure
non-increasing in the agent count. Adding heterogeneous agents
shrinks the residue without eliminating it (composite floor
$\Sfloor(E)>0$).
\end{corollary}

\begin{remark}\label{rem:godel-vs-residue}
\Cref{thm:resid-floor} clarifies the relationship between Gödel's
classical incompleteness theorems and the residue notion: Gödel
establishes the existence of formally undecidable propositions within
fixed formal systems; the residue is the receiver-relative analogue
expressed via $\beta$. Both are positive lower bounds on
representational granularity, but the residue is finer-grained
(continuous in $\beta$ rather than discrete) and admits multi-agent
reduction.
\end{remark}

%==========================================================
\section{Bias-as-decoder, observation-as-projection}
\label{sec:bias}
%==========================================================

\subsection{The classical equivalence}

The classical observation--bias equivalence
(in epistemology, philosophy of perception, and the social-construction
literature) says: functional observation by a finite agent requires
selection criteria; selection criteria constitute systematic bias;
bias and observation are mathematically equivalent.

\subsection{Reformulation}

\begin{theorem}[Bias is decoder, observation is projection]\label{thm:bias-decoder}
For an agent $\agent=(\receiver,\Method,\Goal)$, the classical
``selection criteria'' are precisely the decoder $\D:\X\to\K$ and
the candidate-projection $\Pi:\K\to 2^\X$.
``Functional observation'' is the composite map $\Pi\circ\D$ acting
on $\X$. ``Systematic bias'' is the deviation
$x\mapsto d(x,\Pi(\D(x)))\in[0,\beta]$ from input to canonical
candidate. The equivalence
``observation $\equiv$ bias'' is the trivial identity
``$\Pi\circ\D$ is what the agent does'' rendered as a
floor-positivity statement.
\end{theorem}

\begin{proof}
Selection criteria as $\D,\Pi$: direct from \Cref{def:receiver}. The
``bias'' is the floor $\beta\ge\sup_x d(x,\Pi(\D(x)))$,
strictly positive by hypothesis. The equivalence is a definitional
unpacking once the receiver structure is identified.
\end{proof}

\begin{corollary}[No bias-elimination without receiver-elimination]
\label{cor:no-elim}
Eliminating $\beta$ requires either eliminating the receiver
(no observation) or making the receiver infinite-resolution
($\beta=0$, forbidden by \Cref{thm:floor}). Hence ``observation
without bias'' is mathematically equivalent to ``no observation.''
\end{corollary}

\begin{remark}[Disposing of the equivalence rhetoric]\label{rem:obs-rhetoric}
Once observation and bias are identified with the decoder--projection
composition, the classical claim ``observation is bias'' is no longer
a deep equivalence requiring philosophical
argument; it is a definitional unpacking. The substantive question
is the size of $\beta$, which is decreasable by ensemble composition
(\Cref{thm:multi-cat}) but never zero.
\end{remark}

%==========================================================
\section{Reality-Independence as cell exteriority}
\label{sec:reality-indep}
%==========================================================

\subsection{The reality-event independence claim}

A common claim in the meaning-impossibility literature is that
``reality events occur independently of any observer's interpretation
processes.'' We extract the mathematically substantive content of this
claim and prove it from the cell-truth framework.

\begin{theorem}[Cell exteriority]\label{thm:exterior}
For any action map $\Action:\X\to\Y$, the action-cell
$\Cell_y=\Action^{-1}(\{y\})$ is a property of $(\X,\Action)$ alone
and does not depend on any receiver $\receiver$. In particular, the
cell exists in outcome space and is available to all receivers
simultaneously without coordination.
\end{theorem}

\begin{proof}
By \Cref{def:cell}, $\Cell_y$ is defined entirely from $\X,\Action$,
without any receiver in the definition. Hence it is exterior to any
particular receiver and equally available to all.
\end{proof}

\begin{remark}[Reality-independence as cell exteriority]\label{rem:re-cell}
\Cref{thm:exterior} provides a precise mathematical form of the
informal ``reality is observer-independent'' claim: the action-cell
that constitutes operational truth lives in outcome space, not in any
receiver's representation space. The cell is the same for all
receivers; what differs is how each receiver \emph{decodes to} it.
\end{remark}

\subsection{Coordination without common-knowledge}

\begin{theorem}[Coordination without common knowledge]
\label{thm:coord-no-ck}
For an ensemble $E$ to coordinate on cell $\Cell$, it is sufficient
(by \Cref{thm:ccc}) that each agent's candidate-projection contains
a state at distance $\le\tau(\Cell)-\beta_i$ from $\Cell$. There is
no requirement of common knowledge of $\Cell$, common knowledge of
common knowledge of $\Cell$, or any iterated common-knowledge
hierarchy~\cite{aumann1976,fagin1995}.
\end{theorem}

\begin{proof}
\Cref{thm:ccc} provides the per-agent attainment condition without any
reference to inter-agent knowledge structures. The cell is exterior
to all receivers (\Cref{thm:exterior}), so each agent reaches it
through its own decoder--projection chain independently of what other
agents know.
\end{proof}

\begin{remark}[Disposing of common-knowledge requirements]
\label{rem:ck-disposal}
The classical Lewis--Schelling--Aumann account of coordination relies
on iterated common knowledge of focal points. Cell-coordination via
\Cref{thm:ccc} dispenses with this hierarchy: the cell itself, as an
object in outcome space, plays the role of the focal point; agents
need only decode-to-it, not know-that-others-know-it.
\end{remark}

\begin{figure*}[!htbp]
    \centering
    \includegraphics[width=\textwidth]{validation/figures/panel_5.png}
    \caption{\textbf{Cell-Meaning, Residue, Bias, and Cell Exteriority.}
    (\textbf{A})~Point-meaning forbidden: histogram of projection diameters across $200$ Monte Carlo states (sampled uniformly from $[-2,2]^{2}$) for receiver with $\beta=0.5$ and projection radius $0.5$. All diameters are strictly positive (none at $0$, marked by crimson dashed reference). Empirical content of Theorem~\ref{thm:no-point}: no bounded receiver carries point-meaning; the candidate-projection $\Pi(\D(x))$ is never a singleton when $\beta>0$.
    (\textbf{B})~Gödelian residue equals receiver floor: measured residue (steelblue circles) versus $\beta$ on the linear $x$-axis ($0.1$ to $3.0$), $30$ values. Theory line $y=x$ (crimson) coincides exactly with measurements. Empirical content of Theorem~\ref{thm:resid-floor}: the classical knowledge--understanding gap is exactly the receiver floor $\beta$. Experiment E40 reports relative error $0.00$.
    (\textbf{C})~Cell exteriority: bar chart comparing the receiver floor $\beta_i$ (crimson, predicted) against the mean $S$-value inside the cell (seagreen, measured) for four receivers with $\beta\in\{0.2,0.5,1.0,2.0\}$ and a common cell of tolerance $\tau=1.0$. Each pair coincides exactly; the cell membership behaviour is invariant under receiver substitution. Empirical content of Theorem~\ref{thm:exterior} (Cell exteriority): the action-cell is a property of the outcome space and action map alone, not of any particular receiver.
    (\textbf{D})~Three-dimensional residue surface: composite-floor residue $\prod_i f_i/\Sigma^{n-1}$ over $(n,f)\in\{1,\dots,7\}\times[0.5,8.0]$, plasma colourmap. The surface decreases geometrically in $n$ (Corollary~\ref{cor:resid-ens}: residue reduction by ensemble) and grows linearly in $f$. Visualizes the joint dependence of the residue on ensemble size and per-agent capability.}
    \label{fig:panel5-meaning}
\end{figure*}

%==========================================================
\part{Validation, Connections, and Conclusion}
%==========================================================

\section{Numerical validation}
\label{sec:validation}

\subsection{Methodology}

The validation suite reports forty-five experiments organized into
nine clusters. Each experiment compares a measured quantity to the
theoretical closed-form prediction and reports relative error against
machine precision.

\begin{itemize}[leftmargin=*]
\item \textbf{C1} (E1--E5): receiver foundations --- floor positivity,
attainment, scaling, layered floor, mode non-privilege.
\item \textbf{C2} (E6--E10): cell-truth --- indistinguishability inside
cell, representational invariance under three encodings, layered
cell-closure.
\item \textbf{C3} (E11--E15): multi-agent algebra --- aggregate floor,
catalytic composition (2-, 3-, 5-, 10-agent), composite floor symmetry.
\item \textbf{C4} (E16--E20): common-cell convergence --- disjoint
representation ensembles attaining shared cell, reachability
monotonicity, reachability volume, lower bound, Pareto-reachability.
\item \textbf{C5} (E21--E25): purpose --- existence, uniqueness, stability
under perturbation, robustness to dropout, distribution.
\item \textbf{C6} (E26--E30): motivation heterogeneity --- composite
floor invariance under goal substitution, goal-quotient invariance,
heterogeneity-driven replication reduction.
\item \textbf{C7} (E31--E35): cell- vs.\ point-meaning --- point-meaning
forbidden by floor, cell-meaning generic, eleven-prerequisite
collapse to $\beta=0$.
\item \textbf{C8} (E36--E40): Gödel-residue and bias --- residue equals
floor, residue reduction by ensemble, bias-decoder identification.
\item \textbf{C9} (E41--E45): cell exteriority and coordination ---
cell independence of receiver, coordination without common-knowledge,
asymptotic floor reduction.
\end{itemize}

\subsection{Summary of results}

\begin{table}[h!]\centering
\begin{tabular}{lcc}
\toprule
\textbf{Cluster (\# expts)} & \textbf{Max rel.\ err.} & \textbf{Verdict}\\\midrule
C1 receiver foundations (E1--E5)     & $0.00$                       & PASS\\
C2 cell-truth (E6--E10)              & $0.00$                       & PASS\\
C3 multi-agent algebra (E11--E15)    & $1.89\times 10^{-16}$        & PASS\\
C4 common-cell convergence (E16--E20)& $0.00$ (boolean)             & PASS\\
C5 purpose (E21--E25)                & $0.00$                       & PASS\\
C6 motivation heterogeneity (E26--E30)& $1.81\times 10^{-16}$       & PASS\\
C7 cell- vs.\ point-meaning (E31--E35)& $0.00$ (boolean)             & PASS\\
C8 Gödel-residue and bias (E36--E40)  & $0.00$                       & PASS\\
C9 cell exteriority (E41--E45)        & $0.00$ (boolean)             & PASS\\\midrule
\textbf{Aggregate (45)}              & $\mathbf{1.89\times 10^{-16}}$ & \textbf{45/45 PASS}\\\bottomrule
\end{tabular}
\caption{Validation summary across $45$ experiments. All verdicts PASS
at machine precision. Maximum aggregate relative error
$1.89\times 10^{-16}$, within IEEE-754 double-precision rounding.}
\label{tab:validation}
\end{table}



%==========================================================
\section{Connections to classical literature}
\label{sec:connections}
%==========================================================

\subsection{Game theory and coordination}

Classical game theory and the Schelling--Lewis--Aumann literature on
coordination~\cite{schelling1960,lewis1969,aumann1976,
nash1951,vonneumannmorgenstern1944,harsanyi1977} grounds coordination
in iterated common knowledge of focal points or in equilibrium
selection. \Cref{thm:coord-no-ck} dispenses with the iterated
hierarchy by placing the focal point in outcome space (the cell)
rather than in agents' representations. The closure properties
(\Cref{thm:coalition}) recover the lattice structure familiar from
coalition theory~\cite{harsanyi1977} without the common-knowledge
prerequisite.

\subsection{Distributed computing}

The Fischer--Lynch--Paterson impossibility theorem~\cite{fischer1985}
and the Byzantine generals problem~\cite{lamport1982} establish
fundamental limits on distributed consensus under faulty agents. Our
framework offers a complementary perspective: composite floor
$\Sfloor(E)$ measures the residual ``disagreement granularity'' of a
multi-agent system, and the catalytic composition law
(\Cref{thm:multi-cat}) gives a quantitative reduction recipe. We do
not prove FLP from scratch; we provide a calculus in which FLP-style
impossibilities show up as positive composite-floor
lower bounds.

\subsection{Belief revision and epistemic logic}

The AGM postulates for belief revision~\cite{alchourron1985} and the
Hintikka--Aumann--Halpern--Moses--Vardi framework for reasoning about
knowledge~\cite{hintikka1962,aumann1999,fagin1995} provide elaborate
machinery for tracking individual beliefs and their interactions.
Our framework treats belief as a derived quantity (\Cref{def:belief})
of the receiver's decoder--projection chain. Belief revision in our
calculus corresponds to receiver substitution at fixed cell
(\Cref{thm:cell-closure} from the foundations paper). The cell-truth
framework subsumes the syntactic AGM machinery while abstracting
away the propositional commitments.

\subsection{Information theory and statistics}

Shannon's mutual-information chain rule~\cite{shannon1948,
coverthomas2006} and the Kullback--Leibler additivity over independent
factors~\cite{kullback1951} mirror our catalytic composition law
(\Cref{thm:cat-fact}): the dimensionless deficit $1-\Sfloor/\Sigma$
plays the role of a probability of success, and independent
composition is multiplicative. Fisher information~\cite{fisher1925}
and statistical hypothesis testing~\cite{lehmann2006,wasserman2004}
provide the standard measure-theoretic backdrop for the dispersion
$\sigma$ in methodologies.

\subsection{Foundational logic}

Gödel's incompleteness theorems~\cite{godel1931},
Tarski's undefinability theorem~\cite{tarski1933}, the halting
problem~\cite{turing1936,church1936}, and Chaitin's information-theoretic
incompleteness~\cite{chaitin1974} each establish a strictly positive
floor on what a bounded formal system can decide. Our receiver floor
$\beta$ is the operational analogue: the same kind of
strictly-positive lower bound, but receiver-relative and
multiplicatively decomposable.

\subsection{Functional analysis}

Banach's contraction principle~\cite{banach1922} provides the
fixed-point machinery for methodological floor (\Cref{thm:floor-fp}).
Koopman operators~\cite{koopman1931} are structural analogues of the
decoder--projection lift. Stone duality~\cite{stone1936} suggests an
alternative algebraic packaging of representations as Boolean
algebras. We use these as analogues, not as premises.

\subsection{Combinatorics and category theory}

The partition encoding of \Cref{rem:three-enc} connects to classical
partition theory~\cite{andrews1976,stanley2011}. The categorical
formulation of purposes (\Cref{rem:cat-purpose}), the goal-quotient
functor (\Cref{con:gq-func}), and the OR-monoid structure
(\Cref{prop:boxtimes}) live in standard categorical foundations
\cite{maclane1998,awodey2010,barrwells1990}.

\subsection{Topology and measure}

Reachability volumes (\Cref{def:rvol}) are measure-theoretic
quantities; the underlying topology is supplied by
$d$~\cite{kelley1955,munkres2000}. Cell exteriority
(\Cref{thm:exterior}) is a topological fact about $\X$ alone.

\subsection{Statistical mechanics as analogue}

Borel--Cantelli~\cite{borel1909,cantelli1917}, Boltzmann's entropy
counting~\cite{boltzmann1877}, and the Landauer
limit~\cite{landauer1961,bennett1973} provide structural
analogues for our asymptotic floor result (\Cref{cor:asymptotic})
and for the dimensional-deficit interpretation of $\beta$. We use
them as analogues; no statistical-mechanical premise is assumed.

%==========================================================
\section{Conclusion}
\label{sec:conclusion}
%==========================================================

We have developed a comprehensive mathematical theory of multi-agent
coordination from a single primitive --- the receiver-relative scalar
$S$ measuring residual semantic distance to an action-cell. The
principal results are:

\begin{enumerate}[label=(\arabic*),leftmargin=*]
\item \emph{Cell-Truth foundations} (\Cref{sec:foundations}): the
receiver calculus, the floor $\beta>0$, the action-cell with positive
tolerance, and representational invariance.
\item \emph{Multi-agent algebra} (\Cref{sec:multi-agent}): the
catalytic composition law for $n$-agent ensembles, reachability
monotonicity (\Cref{cor:monotone}), and asymptotic floor reduction
(\Cref{cor:asymptotic}).
\item \emph{Belief Incompatibility and Common-Cell Convergence}
(\Cref{sec:incompatibility,sec:ccc,sec:reach}): agents with mutually
disjoint representation spaces nevertheless attain the same
action-cell (\Cref{thm:ccc}); reachability volume satisfies a
provable lower bound (\Cref{thm:rlb}).
\item \emph{Purpose} (\Cref{sec:purpose,sec:fixedpoint,sec:stability}):
purpose formalized as a fixed-point cell of the multi-agent composite
floor functional, with proven existence, uniqueness up to reachability
equivalence, stability under perturbation, and an $\omega$-limit
characterization for dynamical agents.
\item \emph{Motivation Heterogeneity} (\Cref{sec:motivation,sec:replication,sec:goal-quotient}):
the composite floor depends only on per-agent floors, not on
goal-content; replication's effectiveness derives from heterogeneity,
not repetition; the goal-quotient yields a clean
categorical structure with the OR-monoid as its codomain.
\item \emph{Reformulation of classical impossibilities}
(\Cref{sec:cell-meaning,sec:eleven,sec:godel,sec:bias,sec:reality-indep}):
the eleven meaning-prerequisites collapse to the single condition
$\beta=0$; the Gödelian residue is exactly $\beta>0$; the
observation--bias equivalence is the decoder--projection
composition; reality-independence is cell exteriority; and
coordination requires no common knowledge.
\end{enumerate}

The calculus is independent, self-contained, and validated to machine
precision across forty-five numerical experiments. It provides a
unified framework for reasoning about multi-agent coordination,
purpose, and the relationship between bounded inquiry and operational
truth without requiring shared knowledge, shared concepts,
shared motivation, or shared methodology. What is shared is the
\emph{action-cell}: a structure in outcome space, not in any
receiver's representation. The cell is the coordination object;
agents need only decode-to-it.

\bibliographystyle{plain}
\bibliography{references}

\end{document}